\corrigesChapitre{1}{Calculer dans les complexes}
\begin{corrige}{1}{Parties réelle et imaginaire}
\begin{questions}
\item $\Re=3$, $\Im=-5$.
\item $6-2\mathrm i$.
\item L’unicité de la forme algébrique donne $x=3$ et $2y=-4$, donc $y=-2$.
\end{questions}
\end{corrige}
\begin{corrige}{2}{Produits}
\begin{questions}
\item $1+2\mathrm i+\mathrm i^2=2\mathrm i$.
\item $8-\mathrm i$.
\item $2026=4\times506+2$, donc $\mathrm i^{2026}=\mathrm i^2=-1$.
\end{questions}
\end{corrige}
\begin{corrige}{3}{Conjuguer}
\begin{questions}
\item $-2-7\mathrm i$.
\item $8$.
\item $(4-3\mathrm i)(4+3\mathrm i)=16+9=25$.
\end{questions}
\end{corrige}
\begin{corrige}{4}{Inverses}
\begin{questions}
\item $(1-\mathrm i)/2$.
\item $(3+4\mathrm i)/25$.
\item $(1+\mathrm i)(1-\mathrm i)/2=(1-\mathrm i^2)/2=1$.
\end{questions}
\end{corrige}
\begin{corrige}{5}{Équation linéaire}
\begin{questions}
\item $z=2/(1+\mathrm i)=1-\mathrm i$.
\item $z=(3+\mathrm i)/2$.
\item Si $b=0$, tout complexe convient ; si $b\ne0$, aucune solution.
\end{questions}
\end{corrige}
\begin{corrige}{6}{Un quotient}
\begin{questions}
\item $(3+\mathrm i)(1+\mathrm i)/2=1+2\mathrm i$.
\item $(1+2\mathrm i)(1-\mathrm i)=3+\mathrm i$.
\item $1-2\mathrm i$ ; il vaut également $(3-\mathrm i)/(1+\mathrm i)$.
\end{questions}
\end{corrige}
\begin{corrige}{7}{Avec le conjugué}
\begin{questions}
\item Poser $z=x+\mathrm iy$ : $3x-\mathrm iy=3+4\mathrm i$, donc $z=1-4\mathrm i$.
\item $2\mathrm iy=6\mathrm i$ donne $y=3$, donc $z=x+3\mathrm i$ avec $x\in\R$.
\item Elle fixe seulement la partie imaginaire et laisse la partie réelle libre.
\end{questions}
\end{corrige}
\begin{corrige}{8}{Binôme}
\begin{questions}
\item $1+3\mathrm i+3\mathrm i^2+\mathrm i^3=-2+2\mathrm i$.
\item Le carré de $2\mathrm i$ vaut $-4$.
\item $(-4)^2=16$.
\end{questions}
\end{corrige}
\begin{corrige}{9}{Réalité d’un quotient}
\begin{questions}
\item $z=[x+1+\mathrm i(1-x)]/2$.
\item $1-x=0$, donc $x=1$, donnant $z=1$.
\item $x+1=0$, donc $x=-1$, donnant $z=\mathrm i$.
\end{questions}
\end{corrige}
\begin{corrige}{10}{Conjugaison et puissance}
\begin{questions}
\item Appliquer la conjugaison du produit à $z\times z$.
\item Au rang zéro, les deux membres valent $1$. Si l’identité vaut pour $n$, conjuguer $z^{n+1}=z^nz$ donne $\bar z^n\bar z=\bar z^{n+1}$.
\item Utiliser $z^{-m}=1/z^m$ et la conjugaison de l’inverse : $\overline{z^{-m}}=1/\bar z^m$.
\end{questions}
\end{corrige}
\begin{corrige}{11}{Une équation à analyser}
\begin{questions}
\item $x^2+y^2=2x+2\mathrm iy$ donne $y=0$ et $x^2=2x$.
\item $x(x-2)=0$, donc $z\in\{0,2\}$.
\item Pour $0$, les deux membres sont nuls ; pour $2$, ils valent $4$.
\end{questions}
\end{corrige}
\begin{corrige}{12}{Une somme binomiale}
\begin{questions}
\item $\sum_{k=0}^6\binom6k\mathrm i^k$.
\item $(1+\mathrm i)^2=2\mathrm i$, donc la puissance six vaut $(2\mathrm i)^3=-8\mathrm i$.
\item La partie réelle est nulle : $1-\binom62+\binom64-\binom66=1-15+15-1=0$.
\end{questions}
\end{corrige}
